Ceres és el planeta nan més petit del Sistema Solar i durant molts anys va ser considerat un asteroide, ja que està situat en el cinturó que hi ha entre Mart i Júpiter. Ceres té un període orbital al voltant del Sol de 4,60 anys, amb una massa de $9,43\times 10^{20}\ \mathrm{kg}$ i un radi de 477~km. Calculeu:

\figura[0.45]{fig1.png}

\begin{apartat}{1}
Quin és el valor de la intensitat de camp gravitatori que Ceres crea a la seva superfície? Quina és la velocitat i l'energia mecànica mínima d'una nau espacial que, sortint de la superfície, escapés totalment de l'atracció gravitatòria del planeta?
\end{apartat}

\begin{apartat}{1}
La distància mitjana entre Ceres i el Sol, tenint en compte que la distància mitjana entre la Terra i el Sol mesura $1,50\times 10^{11}\ \mathrm{m}$ i que el període orbital de la Terra al voltant del Sol és d'un any.
\end{apartat}

\blocetiquetat{Dada:}{$G = 6,67\times 10^{-11}\ \mathrm{N\,m^2\,kg^{-2}}$}
